\section{Experimental Methodology}\label{sec:method}
\paragraph{Evidence classes.} Every result is labelled. \textbf{REPLAY}: 64 recorded outputs of the production evaluator for 8 constructed technical questions with 8 constructed answers each (the evaluator itself is not executed), driven through the real hub; the replay candidate never improves, and follow-up questions are answered with the answer to their root question. \textbf{FAULT-INJECTED}: constructed faults in agent output. \textbf{SOCKET}: a local HTTP server with real timeouts. \textbf{SYNTHETIC}: authored feedback texts. No LLM weights were available, so the LLM service is stubbed as unavailable and the hub's non-LLM fallbacks serve follow-ups and feedback; the frozen PPO checkpoint runs unless a condition disables it. There is no human or live-model evidence in this paper.

\paragraph{Protocol and reproducibility.} The protocol, configuration, runners and analysis scripts were committed and tagged before any reported run; the protocol states metrics, units, predictions with criteria, deviations (three definitional fixes made after smoke runs and disclosed there) and exact commands. Each run records the commit, configuration hash and condition flags; results files are write-once, and tables are generated from them. The tag is local and carries no external timestamp.

\paragraph{XA: comparison and ablations.} Units are sessions: 8 recorded-answer sessions $\times$ 20 seeded question orders = 160 per condition, identical inputs across conditions, conditions interleaved in seeded random order. The 160 sessions reuse the same 64 answers and are therefore dependent; intervals are descriptive. Conditions: the original workflow (B0), without PPO (B7), with canonical guardrails (B0c); the full multi-agent configuration (M1); without the Strategy Agent (M2), the Coaching Agent (M3), arbitration (M4), the grounding check (M5), Strategy$\rightarrow$Coaching communication (M6), authority enforcement (M8); M1 without PPO (M7); M1 with canonical guardrails (M1c). Metrics: completion; \emph{authority} (every committed score equals the recorded evaluator score); \emph{transition validity} (levels in $[1,5]$, RL-phase steps $\le1$); level changes and reversals; agreement with the PPO action; guardrail firings; \emph{guardrail bypass} (final actions the guardrail layer would still override); conflicts; agent calls and valid-proposal rate; per-turn latency (pooled p50/p95). Level changes and reversals describe stability and are \emph{not} a quality measure.

\paragraph{XB: agent faults and ablated safeguards.} 8 sessions $\times$ 5 orders = 40 sessions per cell. Faults are injected in the agents' output at rate 0.3 per call (crash, timeout, malformed, forbidden key, stale version, NaN confidence) or on every call (a valid proposal that opposes the evidence; an ungrounded coaching plan), with the real agent running on all other calls. Variants: all safeguards, and each of authority enforcement, arbitration and grounding removed. A fault is \emph{contained} only if the session completes, committed scores remain authoritative and transitions are valid. Proportions use Clopper--Pearson exact intervals~\cite{cp34}.

\paragraph{XC: threshold decomposition and sensitivity (added after XA and XB).} XA's ``no arbitration'' condition (M4) removes four things at once: the confidence threshold $\tau$, the evidence requirement, the one-step bound and the guardrail layer, so it cannot show which one causes its behaviour. After seeing XA and XB we therefore registered a new experiment (a new protocol and tag, written and committed before any XC run, but designed after seeing XA and XB: it is a prospectively tagged experiment, not a preregistration made before the earlier results were known) that changes \emph{only} $\tau\in\{0,0.5,0.6,0.7,0.8,0.9\}$ on the same 160 replay sessions with every other check in force, and five predictions (C1--C5, Sect.~\ref{sec:results}). XC is a new experiment, not an amendment of XA or XB.

\paragraph{Statistics.} The registered XA and XB paired differences use a percentile bootstrap over sessions (10{,}000 resamples, fixed seed); latency differences use pooled-turn quantiles with sessions as resampling clusters. Because the 160 sessions per condition come from only 8 recorded answer sets, session-level intervals treat dependent units as independent and are too narrow~\cite{fw07}. We therefore add, labelled post hoc, a cluster bootstrap over the 8 recorded-answer sessions for the headline XA contrasts, and use it as the primary interval in XC. We call a difference an improvement only if the analysis supports it, and none of XA's behavioural metrics can support that claim.

\paragraph{Pre-registered predictions.} XA: A1--A3 (completion, authority and transition validity everywhere), A4 (no Strategy Agent reproduces the original trajectory), A5 (ablating coaching, communication, grounding or authority leaves difficulty paths unchanged without faults), A6 (no arbitration $\Rightarrow$ guardrail bypass; full $\Rightarrow$ none), A7 (the agents change the difficulty path in $\ge50\%$ of sessions), A8 (median latency overhead $<50$\,ms), A9 (with PPO unavailable the agents still change the path in $\ge50\%$ of sessions), A10 (canonical guardrails change at least one path). XB: B1--B6 (Sect.~\ref{sec:results}). The outcome of every prediction is reported, including contradictions.
